% abstract.tex

\begin{abstract}
Standard signaling theory explains separation from below: low types cannot afford to mimic high types. This paper studies the opposite failure mode. In a dynamic psychological signaling game, agents choose between an ambitious signal and a pooling-safe signal. The ambitious signal not only conveys ability; it also exposes the agent to future judgments about the accuracy of her own self-model. The intensity of this second-order belief penalty follows an exogenous Markov process. When the probability of entering a high-scrutiny state rises, high types may rationally abandon the ambitious signal before high scrutiny is realized.

The mechanism rests on adverse repricing rather than truth verification. Each type carries a primitive self-model-accuracy level, and higher types have more accurate self-models in the primitive environment. High scrutiny is not modeled as a higher-fidelity audit that simply confirms accurate types. It is a social repricing state: with positive probability, exposed self-presentations are adversely reinterpreted, and the loss is proportional to the Bayesian-surprise distance between the type's prior self-accuracy level and the adverse benchmark. This yields scrutiny-type complementarity: higher types carry larger reputation-at-risk and therefore face a larger marginal penalty when judgment intensity rises.

Anticipatory exposure then reverses the standard willingness ranking over signals. Once the transition probability into high scrutiny crosses a threshold, the high type's incentive constraint for the ambitious signal fails. Moreover, in the pooling candidate, D1 assigns any off-path ambitious deviation to the type for whom the deviation is profitable under the wider set of receiver responses. Because future scrutiny makes ambitious exposure most costly for the high type, the off-path ambitious signal no longer receives the high-type interpretation. The unique pure-strategy D1-consistent psychological equilibrium in the binary-signal class is then pooling on the safe signal.

The result is a belief-driven pooling reversal. Separation unravels from the top: good types climb down. A continuous-type extension predicts that ambitious signals need not simply mean-revert in quality near the threshold. Instead, credible ambition is hollowed out: ordinary winners withdraw, while only extreme high types and poorly calibrated low types remain.
\end{abstract}
